\chapter{Conclusion}
\label{chap:conclusion}

\section{Summary}
\begin{draftblock}
This thesis argues that evaluation choices determine what we can truthfully say
about code and reasoning models. The first part showed that metrics must target
the intended capability: CRUXEval measures execution and input-output
reasoning, Counterfeit Conundrum measures understanding of subtle wrong
programs, and LiveCodeBench controls for temporal contamination. Non-owned
supporting case studies add two further examples: SlopCodeBench measures
iterative degradation, and IneqMath separates final-answer accuracy from proof
soundness.

The second part showed that the score alone is rarely enough. CoT and
fine-tuning change which CRUXEval examples are solved, not just average
accuracy. Counterfeit analyses show that generation difficulty does not predict
understanding of near-miss programs. LiveCodeBench reveals overfitting and
scenario-specific differences hidden by older benchmarks. LINC shows that
neurosymbolic reasoning and chain-of-thought prompting can have similar
accuracy while failing in complementary ways.

The third part showed that measurement shapes the systems built afterward.
Execution metrics encourage execution-trace training. Live benchmarks become
standard tools for reasoning-model evaluation. Non-owned examples such as
VibeCodeBench and LeanDojo show how full web-app evaluation and theorem-proving
infrastructure can redirect field practice. ProofOptimizer shows that optimizing
proof length changes the theorem-proving objective itself.
\end{draftblock}

\section{Future Work}
\begin{draftblock}
The immediate future work is to make these evaluation principles more durable.
Live and contamination-resistant benchmarks need infrastructure that can keep
pace with model releases. Soundness-aware reasoning benchmarks need stronger
judges and better ways to audit judge failures. Iterative coding benchmarks need
quality metrics that correspond more directly to developer maintenance cost.
End-to-end web-app evaluation needs evaluator alignment protocols that remain
reliable as applications and agents become more complex.

The broader research direction is to treat evaluation as part of system design.
For code models, this means measuring not only whether generated software works,
but whether it can be understood, repaired, extended, and trusted. For
mathematical and logical reasoning models, it means measuring not only whether
the final answer is correct, but whether the reasoning process is valid and
useful to humans or proof assistants. For the field as a whole, it means
choosing metrics with the knowledge that those metrics will become training
targets.
\end{draftblock}
